%% kapa: exam=1 round=27 year=2016 period=2 form=A subject=appraisal_law subject_ko=감정평가관계법규 seq=10 number=10
%% source: Q-Net 감정평가사 기출문제 (cst00309) artlSeq=5169402 file=extracted/1008/5169402/제27회 감정평가사 1차 문제/제27회 감정평가사 2교시문제(A형).pdf
%% conversion: kordoc 4.24.51 -> tools/gen1.py (auto); review=proofread-sonnet; flags=none
\begin{question}{10}
국토의 계획 및 이용에 관한 법령상 토지거래계약에 관한 허가구역 내에서 법인 아닌 사인(私人)간에 토지를 매매하는 경우 허가를 요하지 않는 토지의 면적 기준으로 옳지 않은 것은? (단, 국토교통부장관 또는 시ㆍ도지사가 기준면적을 따로 정하여 공고하는 경우를 제외함)
\begin{choices}
  \item 주거지역: 180제곱미터 이하
  \item 공업지역: 660제곱미터 이하
  \item 녹지지역: 100제곱미터 이하
  \item 상업지역: 330제곱미터 이하
  \item 도시지역안에서 용도지역의 지정이 없는 구역: 90제곱미터 이하
\end{choices}
\end{question}
\answer{4}
